\newpage
\section{Sottosistema Firmware}

\subsection{Strumenti di Sviluppo}
\subsubsection*{Ambiente di sviluppo: PlatformIO}

Per lo sviluppo, la compilazione e il deploy del firmware è stato adottato l'ecosistema \textbf{PlatformIO} \cite{platformio_os}. A differenza degli IDE tradizionali, PlatformIO si configura come uno strumento professionale da linea di comando (\textit{CLI}) e un sistema di gestione delle build agnostico rispetto all'hardware. 



L'adozione di questa tecnologia ha offerto diversi vantaggi determinanti per il rigore ingegneristico del progetto:
\begin{itemize}
    \item \textbf{Gestione Granulare delle Dipendenze}: Tramite il \textit{Library Manager} integrato, è possibile dichiarare versioni specifiche delle librerie (es. Adafruit GFX, LittleFS), garantendo che il firmware sia riproducibile e non soggetto a \textit{breaking changes} dovute ad aggiornamenti automatici non controllati.
    \item \textbf{Automazione della Build}: PlatformIO automatizza l'intero workflow di compilazione per l'architettura ARM Cortex-M0+ dell'RP2040, gestendo internamente la toolchain \texttt{arm-none-eabi-gcc} e l'integrazione con il framework \textbf{Arduino-mbed} \cite{pico_arduino_core}.
    \item \textbf{Integrazione con Version Control}: Grazie alla configurazione dichiarativa nel file \texttt{platformio.ini}, l'intero ambiente di sviluppo può essere tracciato tramite sistemi di controllo versione (Git), permettendo a diversi sviluppatori di operare su macchine diverse con la certezza di utilizzare i medesimi flag di compilazione (\texttt{-Os}, \texttt{-DDEBUG}) e partizioni di memoria.
\end{itemize}

\paragraph{Configurazione del progetto: platformio.ini}
Il cuore della configurazione risiede nel file \texttt{platformio.ini}. Questo file permette di definire in modo dichiarativo l'hardware di destinazione, la versione del framework e le librerie necessarie, garantendo la riproducibilità dell'ambiente di build su macchine diverse.

Un aspetto fondamentale per questo progetto è stata la configurazione della partizione di memoria dedicata a \textbf{LittleFS}. Tramite il parametro \texttt{board\_build.filesystem\_size}, è stato possibile riservare una porzione specifica della Flash W25Q16 per l'archiviazione dei log.


\paragraph{Workflow da terminale}
L'utilizzo della \textit{Command Line Interface} (CLI) di PlatformIO ha permesso di velocizzare le operazioni di debug e caricamento.

NOTA: comandi principali in appendice!

L'astrazione fornita da PlatformIO permette di gestire la complessità della \textit{toolchain} ARM GCC in modo trasparente, lasciando che lo sviluppatore si concentri sulla logica dell'applicazione e sulla gestione degli interrupt.

\subsubsection*{Librerie e Dipendenze Software}

\paragraph{Analisi dello Stack Software e Dipendenze}
Per l'implementazione del middleware e dei driver di periferica, sono state selezionate librerie ottimizzate per l'ambiente \textit{bare-metal}, privilegiando l'efficienza nell'uso delle risorse e la stabilità del sistema a lungo termine:

\begin{itemize}
    \item \textbf{Parsing Deterministico NMEA (\texttt{minmea})}: È stata adottata la libreria \texttt{minmea} per il suo approccio \textit{heapless}. Essendo scritta in C puro e priva di allocazioni dinamiche della memoria (\texttt{malloc}), garantisce che il parsing delle sentenze avvenga in tempi costanti e senza rischi di frammentazione della RAM. Questo approccio è essenziale per estrarre i campi \texttt{GGA} e \texttt{RMC} senza introdurre \textit{jitter} nell'elaborazione dei dati GNSS.
    
    \item \textbf{Astrazione Grafica (\texttt{Adafruit GFX} \& \texttt{SSD1306})}: La gestione dell'interfaccia OLED sfrutta la libreria \texttt{Adafruit GFX} come \textit{Hardware Abstraction Layer} (HAL) per le primitive grafiche, mentre il driver \texttt{SSD1306} gestisce il trasferimento del \textit{framebuffer} tramite bus I2C a \qty{400}{kHz}. Tale configurazione permette l'aggiornamento dinamico della telemetria riducendo l'impatto sul carico computazionale del core.

    \item \textbf{Sottosistema di Segnalazione (\texttt{FastLED} \& \texttt{NeoPixel})}: Queste dipendenze gestiscono la logica matematica dello spazio colore (RGB/HSV) e la generazione di gradienti. È importante sottolineare che, mentre le librerie gestiscono il calcolo cromatico, la generazione fisica dei segnali di controllo è delegata al sottosistema PIO (Sez. 4.3), risolvendo il conflitto di timing critico tipico dei LED WS2812B.

    \item \textbf{Persistenza e Resilienza (\texttt{LittleFS})}: Integrato nativamente nel core, il filesystem gestisce la logica di \textit{wear leveling} dinamico sulla Flash esterna. La sua architettura transazionale assicura l'integrità dei dati salvati, prevenendo la corruzione dei log in caso di eventi critici come \textit{brown-out} o interruzioni improvvise dell'alimentazione.
\end{itemize}


\paragraph{Parsing GNSS: minmea}

Per l'estrazione delle coordinate geografiche e del timestamp UTC dalle sentenze NMEA fornite dal modulo SAM-M8Q \cite{sam_m8q_ds}, è stata integrata la libreria minmea \cite{minmea_lib}. Si tratta di un parser scritto in C99, estremamente leggero e adatto a sistemi \textit{resource-constrained} poiché non fa uso di memoria dinamica (\textit{heap}), evitando così problemi di frammentazione della RAM durante sessioni di logging prolungate.



Per elaborare i dati in arrivo dalla UART, il firmware identifica l'intestazione della sentenza e la passa al parser per la conversione in strutture dati fruibili:


L'utilizzo di \texttt{minmea} garantisce che il parsing avvenga in un intervallo temporale deterministico, riducendo la latenza di elaborazione e assicurando che il Core 0 possa gestire simultaneamente gli interrupt della seriale e la logica di controllo senza perdita di pacchetti.

\paragraph{Interfaccia Visuale: SSD1306 e Adafruit GFX}

La comunicazione visiva con l'utente avviene tramite un display OLED gestito via I2C. Per astrarre la complessità del disegno dei caratteri e delle primitive grafiche, si utilizzano le librerie Adafruit GFX \cite{adafruit_gfx} e Adafruit SSD1306 \cite{adafruit_ssd1306}. Il display viene inizializzato definendo l'indirizzo I2C (tipicamente \texttt{0x3C}) e allocando il buffer necessario per la gestione dei pixel: L HO MESSO IN APPENDICE

Per massimizzare l'efficienza, il firmware opera tramite un \textit{framebuffer} in RAM: tutte le operazioni di disegno vengono effettuate localmente e trasmesse al controller \texttt{SSD1306} solo tramite la funzione \texttt{display.display()}. Questo approccio riduce il traffico sul bus I2C e previene fenomeni di sfarfallio (\textit{flickering}) durante l'aggiornamento dei dati meteo o delle coordinate in tempo reale.

\paragraph{Feedback Cromatico: FastLED e Offloading PIO}

Per la gestione dei LED WS2812B è stata scelta la libreria FastLED \cite{fastled_lib}, che permette una manipolazione avanzata dello spazio colore \texttt{HSV} (più intuitivo dell'RGB per creare gradienti). Come discusso nella Sezione 4.3, la generazione fisica dei segnali di timing è delegata alle macchine a stati del \textbf{PIO}, mentre la libreria fornisce l'astrazione per definire i colori: MESSO IN APPENDICE

Il sistema utilizza i LED per fornire un feedback immediato sull'integrità del sistema: un gradiente dal rosso al verde indica la precisione del fix satellitare, mentre brevi impulsi luminosi segnalano l'avvenuto salvataggio di un log sulla memoria Flash, garantendo all'utente il monitoraggio dello stato del dispositivo senza dover consultare il display OLED.

\paragraph{Filesystem: LittleFS}

Per la gestione dell'archiviazione dei log e di parametri da salvare in memoria non volatile (Flash W25Q16 \cite{w25q128_ds}), si rende necessario l'uso di un filesystem. É stato scelto il filesystem LittleFS \cite{fastled_lib} \cite{littlefs_repo}, ottimizzato e pensato per memorie con bassa vita in scrittura (100.000 scritture per settore nel nostro caso), fa un uso minimo della ROM e indipendente dalle dimensioni del FS; inoltre é resiliente a perdite di alimentazione improvvise. La documentazione per l'uso dell'API si trova qui \cite{pico_arduino_core}. 

LISTING IN APPENDICE

Il programma utilizza la memoria flash per salvare l'ultima posizione valida del gps in modo che ve ne sia sempre una di fallback.
Anche i dati meteo ritornati dal programma sul PC via UART vengono salvati in flash, permettendo di mantenerli anche se il PC viene scollegato e la board spenta.




\subsubsection*{Offloading Computazionale tramite PIO}

Una delle sfide critiche nello sviluppo del firmware è stata la gestione della concorrenza tra il protocollo dei LED \textbf{WS2812B} e il flusso dati asincrono del modulo GNSS. I LED NeoPixel richiedono una modulazione degli impulsi con timing estremamente rigorosi (nell'ordine dei nanosecondi). Le implementazioni software tradizionali (\textit{bit-banging}) soddisfano questi requisiti entrando in una sezione critica in cui vengono **disabilitati globalmente gli interrupt**. 

Tale approccio è incompatibile con il data-logging GNSS: la disabilitazione degli interrupt, anche per pochi millisecondi, causa l'overflow del buffer della UART, portando alla perdita di byte fondamentali nelle sentenze NMEA. Per ovviare a questo limite, è stato utilizzato il sottosistema \textbf{PIO} (\textit{Programmable I/O}) dell'RP2040:

\begin{itemize}
    \item \textbf{Parallelismo hardware e Determinismo}: Una macchina a stati dedicata (\textit{State Machine}) esegue un micro-codice assembly specifico per generare le forme d'onda. Questo processo avviene in modo totalmente indipendente dai cicli di clock dei core ARM, garantendo un segnale privo di \textit{jitter}.
    \item \textbf{Comunicazione tramite FIFO}: Il \textit{Core 0} carica i dati cromatici in un buffer hardware di tipo \textbf{FIFO} (\textit{First-In-First-Out}). Una volta trasferito il dato, la CPU torna immediatamente a gestire il parsing NMEA, mentre il PIO svuota il buffer pilotando fisicamente il pin.
\end{itemize}



\paragraph{Sincronizzazione e Protocollo WS2812B}

L'implementazione sfrutta un programma PIO che traduce ogni bit di colore in una sequenza di impulsi \texttt{HIGH}/\texttt{LOW} con rapporto ciclico (\textit{duty cycle}) variabile. Questo permette di delegare interamente la gestione del tempo all'hardware dedicato.

Questa architettura realizza un vero e proprio \textbf{offloading hardware}: il processore non deve più attendere il completamento della trasmissione seriale verso i LED (operazione intrinsecamente lenta e bloccante), ma può dedicarsi esclusivamente al mantenimento dell'integrità dei dati di navigazione provenienti dal modulo \textbf{SAM-M8Q}.

\paragraph{Gestione dei Conflitti di Interrupt}

In un sistema embedded senza PIO, l'arrivo di un byte sulla UART durante la trasmissione verso i LED genererebbe una \textbf{Interrupt Latency} inaccettabile. Se l'interrupt venisse servito, il segnale dei LED si corromperebbe (causando sfarfallio); se venisse ignorato, il byte UART andrebbe perso. 



L'offloading su PIO risolve il dilemma alla radice: poiché il PIO non può essere interrotto da eventi software, il protocollo fisico dei LED rimane stabile, mentre i core ARM restano sempre disponibili per rispondere istantaneamente agli interrupt della UART e dei pulsanti GPIO (\textit{falling edge}), garantendo un sistema realmente \textit{real-time}.

\subsection{Periferiche e Persistenza}

Per la memorizzazione non volatile dei log GNSS e dei parametri di sistema, il progetto integra una memoria Flash NOR \textbf{Winbond W25Q128JV} \cite{w25q128_ds}. Il chip, con una capacità di \qty{128}{Mbit} (\qty{16}{MB}), comunica con l'RP2040 tramite interfaccia \textit{Quad-SPI} (QSPI), garantendo un throughput di dati elevato e latenze ridotte durante le operazioni di scrittura intensiva.

\subsubsection*{Modulo GNSS: u-blox SAM-M8Q}
Il sistema di posizionamento si affida al modulo \textbf{u-blox SAM-M8Q} \cite{sam_m8q_ds}, un ricevitore GNSS (\textit{Global Navigation Satellite System}) di fascia professionale che integra un'antenna \textit{patch} ceramica pre-accordata. 

L'integrazione di questo componente (visibile nello schematico di Fig. \ref{fig:gps_schematic}) offre diversi vantaggi tecnologici fondamentali per il progetto:

\begin{itemize}
    \item \textbf{Ricezione Multi-Costellazione}: Il modulo supporta la ricezione simultanea di tre sistemi GNSS (GPS, Galileo e GLONASS). Questa capacità permette di agganciare un numero superiore di satelliti rispetto ai ricevitori a singola costellazione, garantendo una precisione sub-metrica e un fix della posizione estremamente rapido (\textit{Time To First Fix}) anche in condizioni di scarsa visibilità del cielo.
    \item \textbf{Protocollo di Comunicazione}: Il modulo comunica con il microcontrollore tramite interfaccia UART. Le sentenze \textbf{NMEA 0183} vengono trasmesse serialmente e processate dal firmware per estrarre coordinate, altitudine e timestamp UTC. 
    \item \textbf{Integrità del Segnale}: Come evidenziato nello schematico, il modulo è alimentato tramite la linea a \qty{3.3}{V} filtrata dall'LDO dedicato. La scelta di un'antenna integrata minimizza le perdite di segnale dovute a disadattamenti di impedenza, tipiche delle soluzioni con antenna esterna e connettore SMA.
\end{itemize}

L'output seriale del SAM-M8Q viene gestito in modo asincrono dal \textit{Core 0} dell'RP2040 \cite{rp2040_ds}\footnote{Mentre il \textit{Core 1} è dedicato esclusivamente al calcolo della cinematica e all'aggiornamento del display OLED.}, permettendo di mantenere un flusso di log costante sulla Flash senza interferire con le altre funzioni del sistema.

\subsubsection*{File System: Integrità e Wear Leveling con LittleFS}

Data la natura della memoria Flash, soggetta a un numero finito di cicli di cancellazione (tipicamente $10^5$ cicli per settore), la gestione dei dati non avviene tramite scrittura grezza (\textit{raw}), ma attraverso il file system \textbf{LittleFS}. L'adozione di questo file system \textit{open-source}, ottimizzato per microcontrollori, offre vantaggi cruciali per un data-logger professionale:

\begin{itemize}
    \item \textbf{Wear Leveling}: LittleFS distribuisce uniformemente le operazioni di scrittura su tutti i blocchi della memoria, evitando l'usura precoce dei settori iniziali e prolungando drasticamente la vita utile della \textit{W25Q128JV}.
    \item \textbf{Power-loss Resilience}: Il file system è progettato per gestire interruzioni improvvise dell'alimentazione. Grazie a una struttura a "copy-on-write", i dati preesistenti non vengono sovrascritti finché la nuova operazione non è completata, garantendo l'integrità dei log anche in caso di distacco della batteria.
    \item \textbf{Gestione dei Bad Blocks}: Il sistema è in grado di rilevare e isolare dinamicamente eventuali blocchi danneggiati della Flash, mantenendo l'affidabilità del logging nel tempo.
\end{itemize}

Referenza flash figura \ref{fig:flash_schematic}

\subsubsection*{Strategia di Logging e Concorrenza}

Nel firmware sviluppato, la scrittura su LittleFS è affidata esclusivamente al \textit{Core 0}. Per ottimizzare le prestazioni e ridurre i tempi di accesso alla memoria, i dati provenienti dal modulo GNSS \cite{sam_m8q_ds} vengono prima accumulati in un buffer circolare nella RAM dell'RP2040 e successivamente riversati in un file binario giornaliero (es. \texttt{log\_240126.bin}).

Questa separazione logica evita che le operazioni di \textit{I/O} sulla Flash, che possono richiedere diversi millisecondi per la cancellazione di un settore, interferiscano con il campionamento ad alta frequenza dei sensori ambientali gestito dal secondo core.
